\documentclass[12pt]{article}
\usepackage{amsmath, amssymb, geometry, enumitem}
\geometry{margin=1in}
\setlength{\parindent}{0pt}
\setlength{\parskip}{0.8em}

\title{ECON 101: Macroeconomics \vspace{10pt}\\ MCQs for Final Exam Review}
\author{Instructor: Muhebullah Karimzada}
\date{December 03,  2025}

\begin{document}
\maketitle

\begin{enumerate}[label=\textbf{\arabic*.}]

% 1
\item Scarcity exists because:
\begin{enumerate}[label=(\Alph*)]
    \item Markets always fail.
    \item Resources are limited and wants are unlimited.
    \item There are unlimited resources and unlimited wants.
    \item Prices never adjust.
\end{enumerate}

% 2
\item A positive economic statement is:
\begin{enumerate}[label=(\Alph*)]
    \item One that expresses a value judgment.
    \item One that cannot be tested.
    \item One that can be tested and verified with data.
    \item One that must always be true.
\end{enumerate}

% 3
\item The opportunity cost of a choice is:
\begin{enumerate}[label=(\Alph*)]
    \item The sum of all alternatives forgone.
    \item The money cost only.
    \item The highest-valued alternative given up.
    \item The lowest-valued alternative given up.
\end{enumerate}

% 4
\item Comparative advantage occurs when a producer:
\begin{enumerate}[label=(\Alph*)]
    \item Uses more labour than capital.
    \item Produces more with the same inputs.
    \item Has the lowest opportunity cost in producing a good.
    \item Has more resources than competitors.
\end{enumerate}
\newpage
% 5
\item The PPF shifts outward when:
\begin{enumerate}[label=(\Alph*)]
    \item Unemployment increases.
    \item Technology improves.
    \item Prices fall.
    \item A country imports more goods.
\end{enumerate}

% 6
\item Markets are efficient because:
\begin{enumerate}[label=(\Alph*)]
    \item Government sets all prices.
    \item Buyers and sellers respond to incentives.
    \item Scarcity is eliminated.
    \item All goods are public goods.
\end{enumerate}

% 7
\item An increase in the price of a substitute good will:
\begin{enumerate}[label=(\Alph*)]
    \item Reduce demand for the original good.
    \item Increase demand for the original good.
    \item Reduce supply.
    \item Increase supply.
\end{enumerate}

% 8
\item At equilibrium:
\begin{enumerate}[label=(\Alph*)]
    \item There are shortages.
    \item Quantity demanded equals quantity supplied.
    \item Prices rise indefinitely.
    \item There are surpluses.
\end{enumerate}

% 9
\item A decrease in supply with no change in demand will:
\begin{enumerate}[label=(\Alph*)]
    \item Lower the equilibrium price.
    \item Raise the equilibrium price.
    \item Leave price unchanged.
    \item Increase the equilibrium quantity.
\end{enumerate}

% 10
\item A price ceiling below equilibrium causes:
\begin{enumerate}[label=(\Alph*)]
    \item A surplus.
    \item A shortage.
    \item No effect.
    \item Deadweight gain.
\end{enumerate}
\newpage
% 11
\item Deadweight loss is:
\begin{enumerate}[label=(\Alph*)]
    \item Lost revenue for producers.
    \item A reduction in economic efficiency.
    \item An increase in government revenue.
    \item Profit from market power.
\end{enumerate}

% 12
\item A tax on sellers:
\begin{enumerate}[label=(\Alph*)]
    \item Shifts supply left.
    \item Shifts demand right.
    \item Has no effect.
    \item Creates a price floor.
\end{enumerate}

% 13
\item GDP measures:
\begin{enumerate}[label=(\Alph*)]
    \item Total welfare.
    \item Total market value of final goods and services.
    \item Intermediate goods only.
    \item Unpaid household work.
\end{enumerate}

% 14
\item Nominal GDP differs from real GDP because nominal GDP:
\begin{enumerate}[label=(\Alph*)]
    \item Uses base-year prices.
    \item Uses current prices.
    \item Excludes services.
    \item Adjusts for inflation.
\end{enumerate}

% 15
\item If real GDP increases, then:
\begin{enumerate}[label=(\Alph*)]
    \item Production has increased.
    \item Prices have increased.
    \item Inflation has increased.
    \item Imports have decreased.
\end{enumerate}

% 16
\item Long-run growth is driven mainly by:
\begin{enumerate}[label=(\Alph*)]
    \item Money supply growth.
    \item Technological change.
    \item Population growth alone.
    \item Budget deficits.
\end{enumerate}
\newpage

% 17
\item Labour productivity rises when:
\begin{enumerate}[label=(\Alph*)]
    \item Capital per worker increases.
    \item Wages fall.
    \item Unemployment increases.
    \item Trade decreases.
\end{enumerate}

% 18
\item Convergence occurs when:
\begin{enumerate}[label=(\Alph*)]
    \item Corruption increases.
    \item Poor countries adopt strong institutions and policies.
    \item Savings rates fall.
    \item Investment decreases.
\end{enumerate}

% 19
\item In a closed economy:
\begin{enumerate}[label=(\Alph*)]
    \item $S = I$
    \item $S = G$
    \item $I = NX$
    \item $I = C$
\end{enumerate}

% 20
\item The interest rate is determined by:
\begin{enumerate}[label=(\Alph*)]
    \item Money supply only.
    \item Supply and demand for loanable funds.
    \item The government.
    \item Trade balance.
\end{enumerate}

% 21
\item A government budget deficit:
\begin{enumerate}[label=(\Alph*)]
    \item Shifts supply of loanable funds right.
    \item Lowers interest rates.
    \item Shifts supply of loanable funds left.
    \item Increases private saving only.
\end{enumerate}

% 22
\item The MPC is:
\begin{enumerate}[label=(\Alph*)]
    \item $\Delta C / \Delta Y$
    \item $\Delta S / \Delta Y$
    \item $C - Y$
    \item $Y / C$
\end{enumerate}
\newpage
% 23
\item The simple multiplier equals:
\begin{enumerate}[label=(\Alph*)]
    \item $\frac{1}{1 - MPC}$
    \item $1 - MPC$
    \item $MPC - 1$
    \item $1 / MPC$
\end{enumerate}

% 24
\item A fall in autonomous consumption:
\begin{enumerate}[label=(\Alph*)]
    \item Shifts AE upward.
    \item Shifts AE downward.
    \item Raises MPC.
    \item Raises equilibrium GDP.
\end{enumerate}

% 25
\item AD slopes downward because of:
\begin{enumerate}[label=(\Alph*)]
    \item Wealth, interest rate, and international trade effects.
    \item Money illusion.
    \item Sticky wages.
    \item Rising productivity.
\end{enumerate}

% 26
\item LRAS shifts right when:
\begin{enumerate}[label=(\Alph*)]
    \item Prices rise.
    \item Potential GDP increases.
    \item Money supply rises.
    \item AD shifts right.
\end{enumerate}

% 27
\item A negative supply shock:
\begin{enumerate}[label=(\Alph*)]
    \item Shifts SRAS left.
    \item Shifts AD right.
    \item Lowers price level.
    \item Shifts LRAS right.
\end{enumerate}

% 28
\item Money serves as:
\begin{enumerate}[label=(\Alph*)]
    \item A medium of exchange.
    \item A public good.
    \item A capital good.
    \item A consumption good.
\end{enumerate}
\newpage
% 29
\item The simple deposit multiplier is:
\begin{enumerate}[label=(\Alph*)]
    \item $r_d$
    \item $1 - r_d$
    \item $\frac{1}{r_d}$
    \item $r_d - 1$
\end{enumerate}

% 30
\item The Bank of Canada's main policy rate is:
\begin{enumerate}[label=(\Alph*)]
    \item The Bank Rate
    \item The deposit rate
    \item The target overnight rate
    \item The prime rate
\end{enumerate}

% 31
\item Contractionary monetary policy occurs when the Bank of Canada:
\begin{enumerate}[label=(\Alph*)]
    \item Lowers the overnight rate
    \item Raises the overnight rate
    \item Buys government bonds
    \item Depreciates the currency
\end{enumerate}

% 32
\item Higher interest rates tend to:
\begin{enumerate}[label=(\Alph*)]
    \item Increase investment
    \item Reduce investment
    \item Increase net exports
    \item Raise inflation
\end{enumerate}

% 33
\item An appreciation of the Canadian dollar:
\begin{enumerate}[label=(\Alph*)]
    \item Increases net exports
    \item Reduces net exports
    \item Raises exports
    \item Has no effect on trade
\end{enumerate}

% 34
\item In the long run, inflation is caused by:
\begin{enumerate}[label=(\Alph*)]
    \item Productivity growth
    \item Money supply growing faster than output
    \item Tax increases
    \item Trade deficits
\end{enumerate}
\newpage
% 35
\item Unexpected inflation harms:
\begin{enumerate}[label=(\Alph*)]
    \item Borrowers
    \item Lenders
    \item Firms holding nominal assets
    \item Both (B) and (C)
\end{enumerate}

% 36
\item The quantity equation is:
\begin{enumerate}[label=(\Alph*)]
    \item $MV = PY$
    \item $M = PYV$
    \item $P = MVY$
    \item $MV = Y$
\end{enumerate}

\end{enumerate}

\hrule
\vspace{1em}

\begin{center}
\textit{End of Exam MCQs}
\end{center}

\end{document}
